\section{Problem and study contribution}
Robotic motion monitoring needs to distinguish a commanded change from an unauthorized deviation. It also needs a defensible rule for ending an alarm. A predictor can track disturbed measurements well enough to produce a small residual while the process remains displaced from the approved trajectory. Conversely, a nominal reference can retain evidence of that displacement but disagree with legitimate motion when its plant model is inaccurate. This study examines the resulting detection--recovery trade-off under an explicit command-trust assumption.

Our research question is whether a separate return criterion can retain evidence until physical recovery while reducing alarms during legitimate command changes under model mismatch. We test a fixed-return rule selected from earlier development simulations. The distinction between entry and return is operational: entry uses normalized prediction discrepancy; return uses proximity to an approved-command reference, rather than a fraction of the alarm-entry threshold. Neither distinction by itself establishes novelty.

The contribution of the present experiment is a falsifiable test of that mechanism, including controls that receive the same protected context, component removals at both normal-calibrated and identical entry thresholds, and a fresh condition matrix frozen before its outcomes existed. The resulting evidence identifies which component changes recovery and where the overall improvement fails. It does not support a claim that the candidate is a superior detector or an original replacement for existing observer-based diagnosis. The longer-term original-method contribution remains a research objective, rather than a conclusion of this draft.

\section{Related work and comparator scope}
CASPER uses external inertial measurements to monitor robotic motion and studies anomalies introduced through joint-velocity changes \cite{casper}. Its corrected version of record is an ACM Transactions on Internet of Things article; the correction date is distinct from the original 2024 publication. The later edge-computing study investigates real-time robotic-arm anomaly detection \cite{edge}. These papers motivate the concern that a learned predictor and its scoring history may respond differently to brief and persistent deviations. They do not supply an interchangeable benchmark for our four-channel joint-state simulation. Previously saved nine-channel IMU CNNs belong to a separate recorded-data track; no input conversion or claim of paper-faithful reproduction is made here.

TWINE integrates a manufacturing digital-twin framework with explainable anomaly detection \cite{twine}. Its accepted manuscript is available in the author's institutional repository. The inspected supervised classification workflow has different label access from the normal-only procedure used here. Ayankoso et al. combine trajectory-change recognition with condition monitoring \cite{ayankoso}; this is relevant to legitimate motion variation, but does not establish that a change in observed trajectory is authorized. These differences require explicit admission decisions rather than a ranking assembled from published scores.

The methodological background is broader than recent robotic cybersecurity papers. Interval-observer diagnosis already treats the interaction of observer gain, uncertainty and fault sensitivity \cite{meseguer}. CUSUM is an established accumulated-deviation control \cite{nist}. Our three-model envelope, dwell counter and residual fusion therefore receive ordinary-control status. A sampled envelope is not a verified enclosure of all admissible dynamics.

Two recent leads further restrict novelty claims. Ahmed et al. describe a robust adaptive propagated interval observer for actuator-fault diagnosis \cite{rapio}; publisher-indexed information was accessible, but complete recovery equations and matched code remain unverified. Kaisar et al. describe digital-twin-based false-data-injection detection and isolation in robotic vehicles \cite{kaisar}; institutional abstract and bibliographic metadata were verified, but the full manuscript was unavailable to this audit. We do not infer exact equivalence or superiority from these partial records. Their overlap with uncertainty handling and recovery is sufficient to prevent a first-of-its-kind claim until the full comparison is resolved.

The experimental roster consequently consists of transparent controls and our specified candidate, not three faithful recent-paper reproductions. Public voraus-AD recordings are a promising separate source of robot fault observations \cite{voraus}. Their verified availability does not establish an authenticated approved-command channel or the physical-return annotation required by this particular question. No independent recording is silently substituted into the present confirmation.
